% Page 010

\seite{10}

Wie so vieles andere, was die Zeit hervorgebracht, auch der\ob
Zeit zum Opfer gefallen ist, so auch hier, dieses Schloß\ob
lebt nur in dem Denkmal fort, was es sich durch\ob
seine Schenkung in der Pfarrkirche gesetzt hat.

Einen weitern Beleg für das hohe\ob
\margmark{Hohe Persönlich-\\keiten aus\\Urweiche.\\Abt aus Urw.\\s. Joh. Bertels\\„Historia Luxem-\\burgensis\\auch: Geschichte der\\Stadt und der Abtei\\Echternach\\v. Joh. Pet. Brimmeyr, p\\(Imprimerie Centrale\\Gustave Soupert\\Luxemburg. II. Bd.\\1923\\[darüber gekritzelt:]\\Anmerkung\\v. 10.2.1966\\General Kleiber\\aus Urweiche}%
Alter Sefferweichs ist darin zu suchen, daß schon \darueber{ja! 1465}\ob
in dem dreizehnten Jhdt. zwei Brüder nach einander\ob
im Kloster zu Echternach als Äbte segensreich\ob
gewirkt haben, welche ihre Wiege in dem heute\ob
noch stehenden „Kievers“-Hause (Urweich) stehen\ob
hatten. Aber nicht bloß Krummstab verstanden\ob
die Urweicher wohl zu führen, auch mit den Künsten\ob
des Schwertes und Krieges scheinen sie vertraut\ob
gewesen zu sein; denn etwas später wird in\ob
r.\ob
kaiserlichen Diensten ein General Kleiber ge\ohb
nannt, welcher dem, oben genannten „Kievers“-\ob
Hause sein Dasein verdankt. Er starb ohne Erben\ob
und seine Gebeine sollen in Prag ruhen. In seinen\ob
letzten Lebensjahren ersuchte er seine Brüder, ihn\ob
in Prag zu besuchen, aber Prag war weit, und\ob
auf der „Firdel“ hinterm Ofen lag es sich viel un\ohb
gefährlicher, als damals eine Reise durch Deutschland\ob
nach Prag zu machen. Der Besuch unterblieb.\ob
Nach dem Tode des Generals wurden dessen Gebeine \darueber{(kurzum Hab)}\ob
Hinterlassenschaft in Kisten verpackt und in seine\ob
Heimat Urweich gesandt. Die Kisten sollen bloß
\pend
